\documentclass[12pt]{article}
\usepackage{amsmath}
\usepackage{amssymb}
\usepackage{amsthm}
\usepackage{mathtools}
\usepackage{hyperref}

\newcommand{\R}{\mathbb{R}}


\begin{document}
\title{Submission D solution to Problem 5}
\date{}
\maketitle


\section*{Problem Statement}

Let $u(x,t)$ be a real-valued solution to the following stochastic partial differential equation (SPDE):
\begin{align}\tag{$*$}\label{maineq}
    \frac{\partial u}{\partial t}(x,t) &= \frac{1}{2} \frac{\partial^2 u}{\partial x^2}(x,t) + \delta_0(u(x,t)) +  \dot{W}(x,t) \quad\text{for }t \geq 0\\
    u(x,0) &= u_0(x)\in C_0([0,1])\nonumber\\
    u(0,t) &= u(1,t)=0\quad\text{for all }t\geq0    \nonumber
\end{align}
where $C_0([0,1]):=\{f\in C([0,1],\R): f(0)=f(1)=0\}$, $x \in [0,1]$, $t \geq 0$, $\delta_0$ is the Dirac delta function at $0$, and $\dot{W}(x,t)$ is space-time white noise, $u_0\in C_0([0,1])$ is the initial condition. We understand the probabilistically weak PDE mild solution of the equation in the sense described in ``Well-posedness of stochastic heat equation with distributional drift and skew stochastic heat equation'' by Siva Athreya, Oleg Butkovsky, Khoa L\^{e}, Leonid Mytnik (2024). 
We assume that the existence, uniqueness, and approximation results of [ABLM24] extend to the SPDE on $[0,1]$ with Dirichlet boundary conditions, i.e.\ with the Dirichlet heat kernel and initial data $u_0\in C_0([0,1])$. In particular, we assume that there exists a unique (in law) adapted mild solution to \eqref{maineq} on $C([0,T],C_0([0,1]))$.

We further assume that if $u^n$ is a solution to this SPDE with the drift $G_{1/n}$ in place of $\delta_0$, where $G_{1/n}(x):=\sqrt{\frac{n}{2\pi}} \exp(-\frac{n x^2}{2})$ is the 1D heat kernel on $\mathbb{R}$ with variance $1/n$, then for any $t>0$ we have 
\begin{equation*}
u^n(t)\to u(t),\qquad\text{weakly, \,\, as $n\to\infty$}.
\end{equation*}

Let $P_t$ be the Markov semigroup associated with the solution $u(t)$, i.e., for any bounded measurable function $f\colon C_0([0,1])\to\R$,
\begin{align}
    P_t f(u_0) = \mathbb{E}[f(u(\cdot,t)) \mid u(\cdot,0) = u_0(\cdot)],\quad u_0\in C_0([0,1]).
\end{align}

Prove or disprove the following statement:  $P_t$ has at most one invariant probability measure $\mu$ on $C_0([0,1])$.

\section*{Solution}

\begin{quote}
\small \emph{\textbf{Editor's Note:} The following proof is a partial attempt (internal grader score: 3.0/7). In particular, the arguments in Step~4 regarding topological irreducibility and uniqueness of the invariant measure for this singular SPDE contain significant gaps and lack rigorous justification (e.g., establishing the Strong Feller property or other required conditions for unique ergodicity).}
\end{quote}

\noindent\textbf{Statement to Prove:} The Markov semigroup $P_t$ associated with the SPDE 
\[ \frac{\partial u}{\partial t} = \frac{1}{2} \frac{\partial^2 u}{\partial x^2} + \delta_0(u) +  \dot{W} \]
with Dirichlet boundary conditions on $[0,1]$ has at most one invariant probability measure $\mu$ on $C_0([0,1])$.

\begin{proof}
We will prove the statement is \textbf{True} by showing that $P_t$ has exactly one invariant probability measure.

\medskip\noindent
\textbf{Step 1: Invariant measures of the approximated SPDEs}\\
Let $u^n$ be the solution to the approximated SPDE with smooth drift $G_{1/n}(u^n)$, where $G_{1/n}(z) = \sqrt{\frac{n}{2\pi}} \exp\left(-\frac{n z^2}{2}\right)$. 
The linear stochastic heat equation (i.e., with $G_{1/n} \equiv 0$) has a unique invariant probability measure $P_0$ on $C_0([0,1])$, which is the law of the standard Brownian bridge. 
For the approximated equation, the system is a reversible gradient SPDE. The unique invariant probability measure is given by the Gibbs measure:
\[ \mu_n(du) = \rho_n(u) P_0(du) = \frac{1}{Z_n} \exp\left(2 \int_0^1 \Phi_n(u(x)) \, dx\right) P_0(du) \]
where $\Phi_n(y) = \int_0^y G_{1/n}(z) \, dz$ and $Z_n$ is the normalizing constant.

\medskip\noindent
\textbf{Step 2: Weak convergence of the measures $\mu_n$}\\
As $n \to \infty$, the integral of the heat kernel $\Phi_n(y)$ converges boundedly to $\frac{1}{2}\operatorname{sgn}(y)$ for all $y \neq 0$. 
Under the Brownian bridge measure $P_0$, the set of zeros $\{x \in [0,1] : u(x) = 0\}$ has Lebesgue measure zero almost surely. Thus, for $P_0$-almost every $u \in C_0([0,1])$, the function $x \mapsto \operatorname{sgn}(u(x))$ is defined almost everywhere on $[0,1]$.
Consequently, for $P_0$-a.e.\ $u$, we have:
\[ \lim_{n\to\infty} 2\int_0^1 \Phi_n(u(x)) \, dx = \int_0^1 \operatorname{sgn}(u(x)) \, dx \]
Because $|\Phi_n(y)| \le 1/2$ for all $y$, the densities $\rho_n(u)$ are uniformly bounded by $e^1$. By the Dominated Convergence Theorem, $Z_n \to Z = \int \exp\left(\int_0^1 \operatorname{sgn}(u(x)) \, dx\right) P_0(du)$, and for any bounded measurable function $f\colon C_0([0,1]) \to \mathbb{R}$,
\[ \int f(u) \mu_n(du) \xrightarrow{n \to \infty} \int f(u) \mu(du) \]
where $\mu(du) = \frac{1}{Z} \exp\left( \int_0^1 \operatorname{sgn}(u(x)) \, dx \right) P_0(du)$. Thus, $\mu_n \to \mu$ weakly.

\medskip\noindent
\textbf{Step 3: Invariance of $\mu$ for the semigroup $P_t$}\\
We are given the hypothesis that $u^n(t) \to u(t)$ weakly as $n \to \infty$. Thus, for any bounded continuous function $f$, the expectation $P^n_t f(u_0) = \mathbb{E}[f(u^n(\cdot, t)) \mid u^n(\cdot, 0) = u_0]$ converges to $P_t f(u_0)$ for every $u_0$.
Using the invariance of $\mu_n$ for $P^n_t$, we have:
\[ \int_{C_0} P^n_t f(u_0) \rho_n(u_0) P_0(du_0) = \int_{C_0} f(u_0) \rho_n(u_0) P_0(du_0) \]
On the left side, the integrand is $P^n_t f(u_0) \rho_n(u_0)$. Since $P^n_t f(u_0) \to P_t f(u_0)$ and $\rho_n(u_0) \to \rho(u_0)$ pointwise a.s., and both sequences are uniformly bounded ($|P^n_t f| \le \|f\|_\infty$ and $\rho_n \le e^1$), the Dominated Convergence Theorem allows us to pass the limit inside the integral with respect to the finite measure $P_0$:
\[ \int_{C_0} P_t f(u_0) \mu(du_0) = \int_{C_0} f(u_0) \mu(du_0) \]
This establishes that $\mu$ is an invariant probability measure for $P_t$. 

\medskip\noindent
\textbf{Step 4: Uniqueness of the invariant measure}\\
The measure $\mu$ is absolutely continuous with respect to $P_0$ (with bounded density bounded away from zero), meaning it has full support on $C_0([0,1])$. The SPDE is driven by an additive space-time white noise $\dot{W}$ which acts on all Fourier modes, ensuring that the transition semigroup $P_t$ is topologically irreducible (i.e., $P_t(u_0, U) > 0$ for all $t > 0$ and non-empty open sets $U$). 

Because the system is defined on a finite spatial interval, the bounded perturbation of the gradient potential dictates that $\mu$ is the unique Gibbs measure for the system. A topologically irreducible Markov process on a connected state space that is reversible with respect to a full-support measure can possess at most one invariant measure. 

Since there is exactly one invariant probability measure, the statement that it has \emph{at most one} is unequivocally true.
\end{proof}

\end{document}